\section{Xây dựng mô hình YOLO baseline}

\subsection{Dataset}

Tập dữ liệu sử dụng là \textbf{Construction-PPE} do Ultralytics công bố (khoảng
178~MB), gồm ảnh công trường với nhãn hộp bao cho 11 lớp. Tập được chia sẵn thành
train/validation/test như Bảng~\ref{tab:ppe-dataset}, tỷ lệ xấp xỉ 80/10/10.

\begin{table}[H]
\centering
\caption{Tập dữ liệu Construction-PPE}
\label{tab:ppe-dataset}
\begin{tabularx}{\textwidth}{L{3.2cm}Y}
\toprule
\textbf{Thuộc tính} & \textbf{Giá trị}\\
\midrule
Số ảnh & 1.416 (train 1.132, validation 143, test 141)\\
Số lớp & 11\\
Lớp trang bị & \texttt{helmet}, \texttt{gloves}, \texttt{vest}, \texttt{boots}, \texttt{goggles}\\
Lớp thiếu trang bị & \texttt{no\_helmet}, \texttt{no\_goggle}, \texttt{no\_gloves}, \texttt{no\_boots}\\
Lớp khác & \texttt{Person}, \texttt{none}\\
Định dạng nhãn & YOLO (lớp, tâm $x$, tâm $y$, rộng, cao đã chuẩn hóa)\\
\bottomrule
\end{tabularx}
\end{table}

Ngoài tập công khai, đồ án thu thập thêm dữ liệu thực tế gồm các đoạn video và
ảnh công trường (thư mục \texttt{dataset\_raw}) và một tập PPE tùy chỉnh 9 lớp
xuất từ Roboflow (tập validation 229 ảnh) để đánh giá khả năng tổng quát trên dữ
liệu ngoài phân phối huấn luyện. Nhãn được kiểm tra về chỉ số lớp, hộp vượt biên
và ảnh lỗi trước khi huấn luyện; các khung trích từ cùng một video được giữ trong
cùng một tập để tránh rò rỉ ngữ cảnh giữa train và test.

\subsection{Tiền xử lý và tăng cường dữ liệu}

Ảnh được đưa về $640\times640$ bằng letterbox (giữ tỷ lệ, đệm viền), chuyển
BGR$\rightarrow$RGB và chuẩn hóa về $[0,1]$. Các phép tăng cường dùng bộ tăng
cường của Ultralytics: biến đổi HSV, xoay, tịnh tiến, co giãn, cắt xiên, phối
cảnh, lật dọc/ngang, mosaic, mixup và copy-paste. Thay vì dùng hệ số mặc định,
cường độ của các phép này được tìm tự động bằng SHO cùng các siêu tham số tối ưu
hóa (mục tiếp theo). Mosaic được tắt ở 10 epoch cuối (\texttt{close\_mosaic=10})
để mô hình thích nghi với phân bố ảnh thật. Cấu hình tiền xử lý suy luận
(letterbox, thứ tự kênh, chuẩn hóa) được giữ nguyên khi export và trong AI Runtime.

\subsection{Tìm siêu tham số bằng Sea-Horse Optimizer}

Sea-Horse Optimizer \cite{zhao2023sho} là thuật toán tối ưu bầy đàn mô phỏng
hành vi di chuyển, săn mồi và sinh sản của cá ngựa. Mỗi cá thể là một vector
siêu tham số; ở mỗi thế hệ, cá thể di chuyển theo quỹ đạo xoắn ốc hoặc bước
Lévy quanh cá thể tốt nhất (khám phá), săn mồi với xác suất thành công ngẫu
nhiên (khai thác), rồi nửa quần thể tốt nhất sinh thế hệ con. Đồ án hiện thực SHO
độc lập với Ultralytics (\texttt{optimization/sho.py}) để kiểm thử quy tắc cập
nhật bằng hàm mục tiêu toán học trước khi dùng cho YOLO.

Không gian tìm kiếm gồm 18 chiều (Bảng~\ref{tab:sho-space}). Hàm mục tiêu cần
cực tiểu là $f(\mathbf{x}) = 1 - mAP@50{:}95(\mathbf{x})$, trong đó
$mAP@50{:}95(\mathbf{x})$ đo trên tập validation sau một lần huấn luyện \emph{proxy}
ngắn 3 epoch từ cùng checkpoint khởi tạo. Ngân sách tìm kiếm: quần thể 6 cá thể,
3 thế hệ. Bộ siêu tham số tốt nhất được ghi ra \texttt{optimized\_baseline.yaml}
và dùng để huấn luyện baseline đầy đủ.

\begin{table}[H]
\centering
\caption{Không gian tìm kiếm siêu tham số của SHO}
\label{tab:sho-space}
\begin{tabularx}{\textwidth}{L{3.4cm}C{2.8cm}Y}
\toprule
\textbf{Tham số} & \textbf{Khoảng} & \textbf{Ý nghĩa}\\
\midrule
\texttt{lr0} & $[10^{-5}, 10^{-1}]$ & Learning rate ban đầu\\
\texttt{lrf} & $[0{,}01; 1]$ & Tỷ lệ learning rate cuối/ban đầu\\
\texttt{momentum} & $[0{,}6; 0{,}98]$ & Momentum/$\beta_1$\\
\texttt{weight\_decay} & $[0; 0{,}001]$ & Hệ số suy giảm trọng số\\
\texttt{warmup\_epochs} & $[0; 5]$ & Số epoch khởi động\\
\texttt{warmup\_momentum} & $[0; 0{,}95]$ & Momentum giai đoạn khởi động\\
\texttt{hsv\_h}, \texttt{hsv\_s}, \texttt{hsv\_v} & $[0; 0{,}9]$ & Biến đổi màu, bão hòa, độ sáng\\
\texttt{degrees} & $[0; 45]$ & Góc xoay (độ)\\
\texttt{translate}, \texttt{scale} & $[0; 0{,}9]$ & Tịnh tiến, co giãn\\
\texttt{shear} & $[0; 10]$ & Cắt xiên (độ)\\
\texttt{perspective} & $[0; 0{,}001]$ & Phối cảnh\\
\texttt{flipud} & $[0; 1]$ & Xác suất lật dọc\\
\texttt{mosaic}, \texttt{mixup}, \texttt{copy\_paste} & $[0; 1]$ & Xác suất các phép ghép ảnh\\
\bottomrule
\end{tabularx}
\end{table}

\subsection{Huấn luyện}

Baseline được huấn luyện 100 epoch, batch 16, ảnh $640\times640$, trên GPU
Google Colab; toàn bộ đầu ra được ghi trực tiếp lên Google Drive để không mất khi
phiên Colab bị ngắt. Optimizer và learning rate lấy từ kết quả SHO. Seed được cố
định và \texttt{deterministic=True}. Checkpoint tốt nhất được chọn theo
$mAP@50{:}95$ trên tập validation, không chọn theo tập test.

\subsection{Đánh giá baseline}

\begin{table}[H]
\centering
\caption{Kết quả chất lượng của mô hình baseline trên tập test Construction-PPE}
\label{tab:baseline-quality}
\resizebox{\textwidth}{!}{%
\begin{tabular}{lcccccc}
\toprule
\textbf{Model} & \textbf{Params} & \textbf{GFLOPs} & \textbf{Precision} & \textbf{Recall} & \textbf{mAP@50} &
\textbf{mAP@50:95}\\
\midrule
YOLO26n baseline FP32 & 2,57 M & 6,12 & \cbs{} & \cbs{} & \cbs{} & \cbs{}\\
\bottomrule
\end{tabular}}
\end{table}

Ngoài chỉ số tổng hợp, AP theo lớp và confusion matrix được dùng để nhận diện
lớp yếu. Các lớp ``thiếu trang bị'' (\texttt{no\_gloves}, \texttt{no\_goggle}) có
đối tượng nhỏ và ít mẫu, dự kiến là nhóm nhạy cảm nhất với cắt tỉa và lượng tử
hóa, nên được theo dõi riêng ở các giai đoạn sau.
